\chapter{Introduction}
\label{cap:introduccion}

Cryptoassets trade on centralized exchanges (CEX) and decentralized exchanges (DEX), including automated market makers (AMMs) whose quoted prices depend on pool reserves. Arbitrage can move an AMM price toward an external reference, but the amount and speed of that adjustment depend on the swap rule, costs, available reserves, capital and transaction ordering. These mechanisms are studied in the literature on AMM loss-versus-rebalancing (LVR), arbitrage and miner or maximal extractable value (MEV) \cite{milionis2022lvr,daian2020flash,zhou2021hft}. This thesis isolates a much narrower setting: a zero-fee constant-product pool, one exogenous CEX reference and a bounded arbitrage rule.

The first contribution is analytical. A swap of $c$ times the input reserve changes the log price by $2\log(1+c)$ in the direction of the reference, until the no-arbitrage band is reached. This gives an ideal ticks-to-band calculation without simulation. The computational question is what changes when noise, trend and a one-time liquidity-withdrawal trigger act in a fixed schedule, over a finite event horizon.

The agent-based model is deliberately small and reproducible. It represents one constant-product pool, an exogenous reference and stylized noise, trend, arbitrage and liquidity-provider rules. The withdrawal rule is triggered in every tested adaptive-condition run, so it is treated as a fixed intervention rather than evidence of heterogeneous LP adaptation. Depth is varied while absolute order sizes are held fixed in the main design; a separate sensitivity scales orders with reserves. The experiment therefore shows how conclusions depend on the order-to-reserve ratio and panic rule, not how depth alone determines a real market response.

The reduced-model analysis compares persistence, a scalar AR(1), a two-state affine map and a band-aware map on simulator paths. The two-state map does not outperform the scalar autoregression. The band-aware comparison has a small gain concentrated in one synthetic condition and receives the realized future CEX path, so it is a replay diagnostic rather than a real-time forecast. The results do not support a regime-classification or stability claim.

The central finding with practical research value is that first recovery can be temporary: 31.2\,\% of simulated trajectories that returned to the no-arbitrage band later left it, rising to 98.8\,\% at the shallowest tested depth. This motivates distinguishing first return from sustained tracking in future event studies. The thesis contributes a reproducible mechanism-focused experiment and a clear separation between algebraic implications and rule-dependent simulation results; it does not establish effects in live CEX/DEX markets.

The document is organized as follows. Chapter~\ref{cap:problema} delimits the research question and reviews related work; chapter~\ref{cap:justificacion} explains the scientific and methodological contribution. Chapters~\ref{cap:hipotesis} and~\ref{cap:objetivos} state the analytical proposition, evaluation questions and objectives. Chapter~\ref{cap:metodologia} specifies the model and simulation design, chapter~\ref{cap:resultados} reports analytical and computational findings, and chapter~\ref{cap:discusion} discusses their scope. Chapter~\ref{cap:futuro} sets out the work needed for empirical event replay and broader models; the bibliography and technical annexes close the document.
